\documentclass[11pt, a4paper]{article}

% Gemeinsame Style-Definitionen (TEXINPUTS muss templates/ enthalten — siehe scripts/build_tex.py)
\input{bfw-style.tex}

\title{\vspace*{-1.5cm}\Huge\bfseries\color{bfwPrimaryDark}Session~3: Betrieb, Unternehmen \& Unternehmensziele\\[0.4em]
  \large\color{bfwInkSoft}\textnormal{Von JIKU IT-Solutions bis zu Wirtschaftlichkeit, Produktivität und Rentabilität}}
\author{\textsf{IT Lernfeld 1 --- Das Unternehmen und die eigene Rolle im Betrieb beschreiben}}
\date{}

\begin{document}
\maketitle
\thispagestyle{empty}

\vspace{1em}
\begin{merksatz}[title={\bfseries Worum geht es in dieser Session?}]
Jede Auszubildende und jeder Auszubildende arbeitet in einem Betrieb, der zu einem Unternehmen gehört.
Dieses Handout klärt, was genau ein \emph{Betrieb}, ein \emph{Unternehmen} und eine \emph{Firma}
sind --- und warum diese Begriffe fachlich nicht gleichbedeutend sind. Am Beispiel des
Modellunternehmens \emph{JIKU IT-Solutions GmbH} lernen Sie, wie ein modernes IT-Systemhaus
aufgebaut ist, welche Ziele es verfolgt und wie wirtschaftlicher Erfolg mit Kennzahlen wie
\emph{Wirtschaftlichkeit}, \emph{Produktivität} und \emph{Rentabilität} gemessen wird.
\end{merksatz}

\tableofcontents
\vspace{1em}
\hrule
\vspace{1em}

\section{Das Modellunternehmen JIKU IT-Solutions kennenlernen}

\subsection{Was ist ein IT-Systemhaus?}

Ein IT-Systemhaus agiert als Mittler zwischen IT-Herstellern und Endkunden (Unternehmen, Behörden,
Organisationen). Es entwickelt auf Basis von Bestandsaufnahmen vor Ort und Kundenwünschen individuelle
Lösungen, beschafft die erforderlichen Produkte und richtet diese beim Kunden ein. Auch danach übernimmt
es häufig die laufende Wartung der Systeme.

Ein IT-Systemhaus arbeitet sowohl \textbf{per Fernwartung} (direkte Aufschaltung auf Kundensysteme bei
Fehlermeldungen) als auch \textbf{vor Ort} (für persönliche Beratung, Implementierung neuer Komponenten).
Immer mehr Unternehmen geben ihre gesamte IT-Infrastruktur an ein Systemhaus ab
(\emph{IT-Outsourcing}), da sie keine eigene IT-Abteilung mehr unterhalten wollen.

\subsection{Eckdaten JIKU IT-Solutions}

\begin{center}
\begin{tabular}{p{5.5cm} p{7.5cm}}
\toprule
\textbf{Merkmal} & \textbf{Angabe}\\
\midrule
Systemhäuser im Verbund & 10 Systemhäuser\\
Standorte in Deutschland & 16 Standorte\\
Mitarbeiter (gesamt) & 480 Mitarbeiterinnen und Mitarbeiter\\
Mitarbeiter am Standort Hamburg & 52 Mitarbeiter, davon 8 Auszubildende\\
Auszubildende (gesamt) & 96 Auszubildende\\
Besonderheit & Eigene Rechenzentren; jedes Systemhaus rechtlich selbstständig\\
\bottomrule
\end{tabular}
\end{center}

\begin{merksatz}[title={\bfseries Merke: Struktur des JIKU-Verbunds}]
Rechtlich ist jedes Systemhaus im JIKU-Verbund \textbf{selbstständig}. Zusätzlich sind alle
Systemhäuser mit unterschiedlichen Anteilen an einem \textbf{Gemeinschaftsunternehmen} beteiligt,
das z.\,B.\ zentrale Dienste anbietet. Einzelne Systemhäuser erhalten im Verbund
\textbf{Schwerpunktaufgaben}.
\end{merksatz}

\subsection{Leistungsportfolio und Corporate Social Responsibility}

Das Leistungsangebot von JIKU IT-Solutions umfasst vier Hauptbereiche:

\begin{itemize}
  \item \textbf{IT-Service:} Vor-Ort-Service, IT-Betreuung, IT-Management, IT-Outsourcing, IT-Vertrieb, Help-Desk und User-Help-Desk
  \item \textbf{Cloud-Hosting:} Cloud-Migration, Business-/Server-Hosting, Managed Hosting, Rechenzentrum Housing \& Colocation
  \item \textbf{Server \& Storage:} Optimierte Serversysteme, virtualisierte Server, Virtual Desktop Infrastructure (VDI), hyperkonvergente Serversysteme
  \item \textbf{IT-Infrastrukturen:} Server- und Storage-Systeme, Standortvernetzung (VPN), Telefonanlagen \& IP-Telefonie, Virtualisierung
\end{itemize}

Im Bereich \textbf{Corporate Social Responsibility (CSR)} setzt sich JIKU für nachhaltige Unternehmensführung,
Work-Life-Balance (flexible Arbeitszeiten, Home-Office, Betriebssport) und IT-Nachwuchsförderung
(Girls Day, Hackathon-Sponsoring, Stipendien) ein.

\section{Betriebe und Unternehmen im Umfeld unterscheiden}

\subsection{Begriffsabgrenzung: Betrieb, Unternehmen, Firma}

Die drei Begriffe werden im Alltag oft gleichbedeutend verwendet, sind fachlich aber klar zu trennen:

\begin{center}
\begin{tabular}{p{3cm} p{10cm}}
\toprule
\textbf{Begriff} & \textbf{Definition}\\
\midrule
Unternehmen & Rechtlich-wirtschaftliche Einheit mit dauerhafter Tätigkeit und Gewinnabsicht,
die Bücher führt und eine Steuererklärung abgibt. Ein Unternehmen kann mehrere
Betriebe haben.\\
\addlinespace
Betrieb & Organisatorisch-technische Wirtschaftseinheit, die Güter und Dienstleistungen
beschafft, erstellt, anbietet und verkauft (z.\,B.\ Handelsbetrieb, IT-Betrieb,
Rechenzentrumsbetrieb).\\
\addlinespace
Firma & Nach \textbf{§~17 HGB} der Name, unter dem ein Kaufmann seine Geschäfte betreibt
und seine Unterschrift abgibt. Die Firma ist der \emph{Handelsname} --- nicht das
Unternehmen selbst.\\
\bottomrule
\end{tabular}
\end{center}

\begin{merksatz}[title={\bfseries Merke: Unternehmen und Betrieb}]
Ein \textbf{Unternehmen} kann mehrere \textbf{Betriebe} haben. Beim JIKU-Verbund ist jedes
Systemhaus rechtlich ein eigenes Unternehmen --- die einzelnen Standorte sind dessen Betriebe.
\end{merksatz}

\subsection{Erwerbswirtschaftliche und gemeinwirtschaftliche Unternehmen}

\begin{center}
\begin{tabular}{p{4.5cm} p{8.5cm}}
\toprule
\textbf{Unternehmenstyp} & \textbf{Merkmale}\\
\midrule
Erwerbswirtschaftlich & Privatunternehmen; handeln auf eigene Rechnung und Verantwortung;
streben nach Gewinn (Gewinnmaximierung). Der Gewinn fließt Inhabern,
Gesellschaftern oder Kapitalgebern zu. JIKU IT-Solutions ist ein
erwerbswirtschaftliches Unternehmen.\\
\addlinespace
Gemeinwirtschaftlich & Staatliche oder gemeinnützige Betriebe; Ziel ist Kostendeckung und
Gemeinwohl (keine Gewinnausschüttung). Beispiele: Stadtwerke,
kommunale Theater, Sozialunternehmen.\\
\bottomrule
\end{tabular}
\end{center}

\subsection{Klassifikation von Betrieben nach Wirtschaftszweigen}

\begin{center}
\begin{tabular}{p{3.5cm} p{4.5cm} p{5cm}}
\toprule
\textbf{Sektor} & \textbf{Inhalt} & \textbf{Beispiele}\\
\midrule
Primärer Sektor & Urproduktion & Landwirtschaft, Forstwirtschaft, Rohstoffgewinnung\\
Sekundärer Sektor & Produktion und Verarbeitung & Maschinenbau, Chipfertigung, Elektronikindustrie\\
Tertiärer Sektor & Dienstleistungen und Handel & IT-Systemhäuser, Banken, Versicherungen, Handel\\
\bottomrule
\end{tabular}
\end{center}

In Deutschland gibt es fast \textbf{4 Millionen Unternehmen}, davon weit über 95\,\% als kleine und mittlere
Unternehmen (\emph{KMU}). Die ITK-Branche zählt mit über \textbf{1 Million Arbeitskräften} in den
Unternehmen zu einer der größten Branchen.

\subsection{Wirtschaftliche Verflechtungen}

Unternehmen kooperieren und schließen sich auf verschiedene Weisen zusammen:

\begin{center}
\begin{tabular}{p{3.5cm} p{9.5cm}}
\toprule
\textbf{Form} & \textbf{Beschreibung}\\
\midrule
Konzern & Zusammenschluss größerer, rechtlich selbstständiger Unternehmen zu einer großen
wirtschaftlichen Einheit unter einer gemeinsamen Leitung (Holding)\\
Fusion/Merger & Zusammenschluss rechtlich selbstständiger Unternehmen zu einer
wirtschaftlichen und rechtlichen Einheit\\
Kartell & Vertraglicher Zusammenschluss zur Erlangung von Wettbewerbsvorteilen;
Preis-/Gebietskartelle sind \textbf{verboten} (GWB, Bundeskartellamt)\\
Franchise & Geschäftsmodell: Franchisegeber überlässt rechtlich und finanziell
selbstständigen Franchisenehmern ein Konzept gegen Entgelt\\
Konsortium & Vorübergehender Zusammenschluss zur gemeinsamen Durchführung eines
größeren Projekts oder Standardisierungsvorhabens\\
\bottomrule
\end{tabular}
\end{center}

Zusammenschlüsse können \textbf{horizontal} (gleiche Leistungsstufe), \textbf{vertikal} (mehrere
Leistungsstufen) oder \textbf{diagonal/anorganisch} (ohne Bezug zur Leistungsstufe) erfolgen.

\section{Ziele von Betrieben und Unternehmen}

\subsection{Unternehmensleitbild von JIKU IT-Solutions}

Das Leitbild beschreibt das Selbstverständnis und die Grundsätze eines Unternehmens. Es wird schriftlich
fixiert und richtet sich nach innen (Mitarbeiter) und nach außen (Kunden, Öffentlichkeit):

\begin{center}
\begin{tabular}{p{2.5cm} p{10.5cm}}
\toprule
\textbf{Element} & \textbf{Inhalt beim JIKU-Leitbild}\\
\midrule
Vision & „Wir wollen immer das beste IT-Systemhaus unserer Region sein und diesen Status halten."\\
Mission & „Alle unsere Handlungen werden wirtschaftlich, kompetent, sozial nachhaltig und umweltfreundlich
ausgeführt; unsere Kunden schätzen dies als Alleinstellungsmerkmal mit dem Motto: ,Die können es einfach!'"\\
Werte & Kundenorientierung, Kompetenz/Exzellenz, Wirtschaftlichkeit, Nachhaltigkeit\\
\bottomrule
\end{tabular}
\end{center}

\subsection{Übersicht der Unternehmensziele}

Ziele lassen sich nach verschiedenen Kriterien einteilen:

\begin{center}
\begin{tabular}{p{3.5cm} p{9.5cm}}
\toprule
\textbf{Zieltyp} & \textbf{Beschreibung}\\
\midrule
Sachziel (Leistungsziel) & Bezieht sich auf Art, Menge, Qualität, Ort und Zeit der erzeugten
Güter oder erbrachten Dienstleistungen; z.\,B.\ Bereitstellung von
5\,000 IT-Arbeitsplätzen, Reduktion von Systemausfällen um 20\,\%\\
\addlinespace
Formalziel & Beschreibt den allgemeinen wirtschaftlichen Erfolg; i.\,d.\,R.\ vor den
Sachzielen formuliert; z.\,B.\ Umsatzsteigerung, Gewinnmaximierung,
Marktanteilsgewinn, Produktivitätssteigerung\\
\addlinespace
Strategisches Ziel & Langfristig; von der Unternehmensleitung festgelegt;
z.\,B.\ Technologieführerschaft, Marktführerschaft\\
\addlinespace
Operatives Ziel & Kurz- bis mittelfristig; von Abteilungs- oder Teamleitung festgelegt;
z.\,B.\ Absatzmaximierung, Kostenminimierung\\
\bottomrule
\end{tabular}
\end{center}

\begin{merksatz}[title={\bfseries Zielformulierung nach der SMART-Formel}]
Ziele sollen \textbf{SMART} formuliert sein: \textbf{S}pezifisch (eindeutig, präzise) ---
\textbf{M}essbar --- \textbf{A}usführbar/Akzeptiert (erreichbar, angemessen) ---
\textbf{R}ealistisch (möglich) --- \textbf{T}erminierbar (klare Zeitvorgabe). Nur überprüfbare
Ziele können konsequent verfolgt werden.
\end{merksatz}

\subsection{Stakeholder und ihre Interessen}

Unternehmen müssen auf die Ansprüche verschiedener Interessengruppen eingehen:

\begin{center}
\begin{tabular}{p{3.5cm} p{9.5cm}}
\toprule
\textbf{Stakeholder} & \textbf{Typisches Interesse}\\
\midrule
Mitarbeiter & Hohes Einkommen, Arbeitsplatzsicherheit, guter Arbeitsplatz\\
Kapitalgeber & Hohe Rendite, Kapitalsicherheit, Sonderausschüttungen\\
Lieferanten & Zahlungsfähigkeit des Unternehmens, hoher Umsatz, Vertrauen\\
Kunden & Günstiger Preis, hohe Qualität, guter Service, schnelle Mängelregulierung\\
Staat/Kommunen & Hohe Steuereinnahmen, viele Arbeitsplätze, Einhaltung von Gesetzen\\
Mitbewerber & Fairer Wettbewerb, Kooperationsmöglichkeiten\\
\bottomrule
\end{tabular}
\end{center}

\subsection{Zielbeziehungen}

Ziele können sich gegenseitig unterstützen oder behindern:

\begin{center}
\begin{tabular}{p{4cm} p{4.5cm} p{4.5cm}}
\toprule
\textbf{Zielbeziehung} & \textbf{Beschreibung} & \textbf{Beispiel}\\
\midrule
Komplementär & Ziele unterstützen/verstärken sich & Höhere Energieeffizienz → Kosteneinsparung\\
Konkurrierend & Ziele behindern sich gegenseitig & Mehr Innovationsausgaben ↔ Kostensenkung\\
Indifferent/Neutral & Ziele beeinflussen sich nicht & Neue Betriebskantine \& Serverausbau\\
\bottomrule
\end{tabular}
\end{center}

\subsection{Wirtschaftliche Kennzahlen}

\begin{center}
\begin{tabular}{p{3.5cm} p{6cm} p{4cm}}
\toprule
\textbf{Kennzahl} & \textbf{Bedeutung} & \textbf{Formel}\\
\midrule
Wirtschaftlichkeit & Verhältnis von Erträgen zu Aufwendungen; Wert $>1$ zeigt Effizienz &
$\dfrac{\text{Erträge}}{\text{Aufwendungen}}$\\[1.2em]
Produktivität & Verhältnis von Output zu Input; misst die mengenmäßige Leistungseffizienz &
$\dfrac{\text{Output}}{\text{Input}}$\\[1.2em]
Rentabilität & Verzinsung des eingesetzten Kapitals in Prozent &
$\dfrac{\text{Jahresüberschuss} \times 100}{\text{Kapital}}$\\[0.6em]
\bottomrule
\end{tabular}
\end{center}

\begin{merksatz}[title={\bfseries Rechenbeispiel Wirtschaftlichkeit}]
JIKU erzielt an einem Tag Erträge von 12\,000\,€ bei Aufwendungen von 10\,000\,€.
Wirtschaftlichkeit = 12\,000 $\div$ 10\,000 = \textbf{1{,}2} (bzw.\ 120\,\%).
Das Unternehmen arbeitet wirtschaftlich, da der Wert über 1 liegt.
\end{merksatz}

\end{document}
